\centerline{\bf \Huge Отбор на подготовительные курсы 7-IT}

\begin{flushright}
        \scriptsize{\textbf{\textit{Именно математика даёт надёжнейшие правила: тому кто им следует — тому не опасен обман чувств. \\Леонард Эйлер}}}
\end{flushright}

\hspace{3mm}

\problem{\textbf{1}} Начальник записал на доске следующую сумму равную $X$: 

$2025\cdot2024 + 2023\cdot2022 + 2021\cdot2020 + 2019\cdot2018 + 2017\cdot2016 + 2015\cdot2014= X$

Чему равна последняя цифра числа $X$?

\hspace{3mm}

\centerline{\textbf{\textit{Решение:}}}

\textbf{Первое решение (\textit{красивое}):}

$2025\cdot2024$ данное произведение имеет такую же последнюю цифру, как и произведение последних цифр данных чисел $5\cdot4=20$ => последняя цифра 0.

Аналогично с остальными:

$2023\cdot2022$ => $3\cdot2=6$ => последняя цифра 6.

$2021\cdot2020$ => $1\cdot0=0$ => последняя цифра 0.

$2019\cdot2018$ => $9\cdot8=72$ => последняя цифра 2.

$2017\cdot2016$ => $7\cdot6=42$ => последняя цифра 2.

$2015\cdot2014$ => $5\cdot4=20$ => последняя цифра 0.

Последняя цифра суммы имеет такую же последнюю цифру, как и сумма последних цифр слагаемых:

$0+6+0+2+2+0=10$ => последняя цифра 0.

\textbf{Ответ: 0}

\newpage
\hspace{3mm}

\textbf{Второе решение (\textit{в лоб}):}

Посчитаем все произведения и найдём их сумму:

$2025\cdot2024=4098600$

$2023\cdot2022=4090506$

$2021\cdot2020=4082420$

$2019\cdot2018=4074342$ 

$2017\cdot2016=4066272$ 

$2015\cdot2014=4058210$ 

$4098600 + 4090506 + 4082420 + 4074342 + 4066272 + 4058210 = 24470350$ => последняя цифра 0.

\textbf{Заодно убеждаемся в верности первого решения!}

\textbf{Ответ: 0}





\hspace{3mm}

\problem{\textbf{2}} В отборе участвуют только шестиклассники, рассаженные в пять кабинетов. В первом кабинете шестиклассников в два раза меньше, чем во втором, во втором на одного шестиклассника меньше, чем в третьем, в третьем столько же, сколько и в четвертом, а в четвертом на 50\% меньше, чем в пятом. Может ли в пятом кабинете сидеть шестиклассников ровно 36\% от общего числа шестиклассников, участвующих в отборе?

\hspace{3mm}

\centerline{\textbf{\textit{Решение:}}}

Пусть в первом кабинете $X$ шестиклассников, тогда во втором — $2X$ шестиклассников, так как по условию задания - в первом кабинете шестиклассников в два раза меньше, чем во втором.

В третьем кабинете $2X+1$ шестиклассников, так как по условию задания - во втором на одного шестиклассника меньше, чем в третьем.

\newpage
\hspace{3mm}

В четвёртом кабинете $2X+1$ шестиклассников, так как по условию задания - в третьем столько же, сколько и в четвёртом.

В пятом кабинете $4X+2$ шестиклассников, так как по условию задания - в четвёртом на 50\% меньше, чем в пятом.

Предположим, что в пятом кабинете ровно 36\% от общего числа шестиклассников.

Тогда:

В четвёртом кабинете ровно 36\%:2=18\% от общего числа шестиклассников — 

$(4X+2):2=(2X+1)$.

В третьем кабинете ровно 18\% от общего числа шестиклассников — также $(2X+1)$.

Во втором кабинете меньше 18\% от общего числа шестиклассников — $(2X+1)>2X$.

В первом кабинете меньше 9\% от общего числа шестиклассников — $2X:2=X$.

Имеем суммарно меньше 99\% = 36\% + 18\% + 18\% + 18\% + 9\% 

А должно быть ровно 100\%

Следовательно, наше предположение о том, что в пятом кабинете ровно 36\% от общего числа шестиклассников было неверным.

\textbf{Ответ: Нет, не может.}


\hspace{3mm}

\problem{\textbf{3}} Чему равно следующее выражение?

$\left(0,025\cdot2000 : \left(0,4 + 1\dfrac{3}{5}\right) + \left(6\dfrac{3}{4}\cdot\dfrac{4}{9} + \dfrac{1}{2} : 0,5\right)\cdot\dfrac{0,6}{0,12}\right):$

$:\dfrac{0,62374-0,17374}{\left(89\cdot98-8650-\dfrac{3\cdot3\cdot3\cdot3\cdot3\cdot3\cdot3}{81}\right)}=?$

\newpage
\hspace{3mm}


\centerline{\textbf{\textit{Решение:}}}



\textbf{Решим по действиям:}

1) $\left(0,4 + 1\dfrac{3}{5}\right)=0,4+1,6=2$

2) $0,025\cdot2000 :2=25$

3) $6\dfrac{3}{4}\cdot\dfrac{4}{9} + \dfrac{1}{2} : 0,5=\dfrac{27}{4}\cdot\dfrac{4}{9} + \dfrac{1}{2}:\dfrac{1}{2}=3+1=4$

4) $\dfrac{0,6}{0,12}=\dfrac{60}{12}=5$

5) $4\cdot5=20$

6) $25+20=45$

7) $0,62374-0,17374=0,45$

8) $89\cdot98=8722$

9) $\dfrac{3\cdot3\cdot3\cdot3\cdot3\cdot3\cdot3}{81}=\dfrac{3\cdot3\cdot3\cdot3\cdot3\cdot3\cdot3}{3\cdot3\cdot3\cdot3}=3\cdot3\cdot3=27$

10) $8722-8650-27=45$

11) $45:\dfrac{0,45}{45}=\dfrac{45}{1}\cdot \dfrac{45}{0,45}=45\cdot100=4500$

\textbf{Ответ: 4500}

\hspace{3mm}
\hspace{3mm}

\problem{\textbf{4}} Аркадий и Савр решили устроить лёгкую пробежку вокруг ЭТЛ по часовой стрелке. Каждый раз, когда Аркадий обгоняет Савра, он называет вслух число, равное количеству обгонов им Савра (сразу после первого обгона называет число 1, сразу после второго обгона называет число 2 и так далее). Чему равна сумма всех чисел, названных Аркадием с начала пробежки и до момента, когда он пробежит ровно 2026 кругов вокруг ЭТЛ, если известно, что мальчики начали одновременно с одного места?

\hspace{3mm}

\centerline{\textbf{\textit{Решение:}}}

В отличие от пятого задания, ученики должны были \textbf{самостоятельно} \textbf{понять}, что в условии чего-то не хватает.

\newpage
\hspace{3mm}

По условию Аркадий должен обгонять Савра.


Но во сколько раз Аркадий быстрее Савра?

Именно этого не хватает в условии.

\textbf{Первое решение (не более 3 баллов):}

Условие задания неполное. Где скорости Аркадия и Савра???

\textbf{Второе решение (не менее 3 баллов):}

Условие задания неполное. Пусть скорость Аркадия такая-то, а скорость Савра такая-то \textbf{(ученик волен в своем выборе)}. И далее следует решение, основанное на предположении скоростей \textbf{(необязательно полное)}.

\textbf{НЕСТАНДАРТНАЯ ЗАДАЧА ДЛЯ УЧЕНИКА, ТРЕНИРУЮЩАЯ ЕГО НАВЫКИ КРИТИЧЕСКОГО МЫШЛЕНИЯ <3}



\hspace{3mm}



\problem{\textbf{5}} Среди 2025 учеников ЭлМШ есть \textit{шиншиллы}, которые всегда говорят правду, и \textit{лососи}, которые всегда лгут (есть хотя бы одна \textit{шиншилла} и хотя бы один \textit{лосось}). Какое наибольшее количество \textit{лососей} может быть, если каждый сказал: "Всего среди 2025 учеников не более 26 \textit{шиншилл}"?

Есть ли логическая ошибка в условии задания? Если есть, объясните её и предложите решение этой проблемы. Если нет, решите задание.

\hspace{3mm}

\centerline{\textbf{\textit{Решение:}}}

А вот в пятом задании присутствует явный намёк на проблему в условии задания. Эту проблему нужно подробно объяснить и предложить решение \textbf{(исправление условия)}.

Заметим, что по условию среди учеников должна быть хотя бы одна шиншилла и хотя бы один лосось. При этом они все произносят одну и ту же фразу. \textbf{Чего не может быть!} Так как шиншилла всегда говорит правду, а лосось всегда лжёт \textbf{(одно и то же высказывание не может быть одновременно и правдой, и ложью)}.


\newpage
\hspace{3mm}

\textbf{Одно из множества решений проблемы: }

Убрать из условия задания следующее предложение:

\textbf{(есть хотя бы одна \textit{шиншилла} и хотя бы один \textit{лосось})}

\textbf{Первый вариант:}

Тогда все 2025 учеников могут быть лососями. \textbf{Однако!}

В этом случае все лососи говорят правду, что противоречит условию. Следовательно, требуется \textbf{второе изменение в условии}:

"Всего среди 2025 учеников \textbf{более} 26 шиншилл"

\textit{Все лососи говорят ложь, и наибольшее количество лососей - 2025.}

\textbf{Второй вариант:}

Тогда все 2025 учеников могут быть шиншиллами. \textbf{Однако!}

В этом случае все шиншиллы говорят ложь, что противоречит условию. Поэтому снова требуется \textbf{второе изменение в условии}:

"Всего среди 2025 учеников \textbf{более} 26 шиншилл"

\textit{Все шиншиллы говорят правду, и наибольшее количество лососей - 0.}

\textbf{ЕЩЕ ОДНА НЕСТАНДАРТНАЯ ЗАДАЧА ДЛЯ УЧЕНИКА, ТРЕНИРУЮЩАЯ ЕГО НАВЫКИ КРИТИЧЕСКОГО МЫШЛЕНИЯ <3}


\hspace{3mm}

\problem{\textbf{6}} Однажды Елена решила отправиться из Оргакина (посёлок в Ики-Бурульском районе) в разностороннее путешествие. На своё усмотрение за одно действие она может поехать либо строго на юг, либо строго на север. За первое действие Елена проехала 1км, а за каждое последующее действие в два раза больше предыдущего (2км, 4км, 8км, 16км и так далее). Может ли Елена вернуться в Оргакин за некоторое количество таких действий (Оргакин не объёмный и является точкой на плоскости)?

\newpage
\hspace{3mm}

\centerline{\textbf{\textit{Решение:}}}

Изначально Елена находится в Оргакине, то есть расстояние от неё до Оргакина равно 0 км — это чётное значение.

Затем Елена удаляется на 1км, и расстояние от неё до Оргакина становится равным 1 км — нечётное значение.

Все последующие действия изменяют расстояние от Елены до Оргакина на чётное количество км. Следовательно, они не меняют чётность значения, которое было нечётным после первого действия.

Получается, что расстояние от Елены до Оргакина после первого действия и после всех последующих всегда измеряется нечётным количеством км. Поэтому оно не может стать равным чётному значению, в частности, 0 км. Значит, вернуться в Оргакин она не сможет.

\textbf{Ответ: Нет, не может.}


\hspace{3mm}

\problem{7} Адьян пришел на III ежегодную Математическую Олимпиаду ЭлМШ x ЭТЛ, на которой было предложено 5 задач. Адьян очень торопился и мог потратить на Олимпиаду максимум час. Задачи, придуманные Расулом Владимировичем, Адьян читает за 2 минуты, а решает за 35 минут. Задачи, придуманные Анджуром Аркадьевичем, он читает 2 минуты, а решает за 15 минут. Задачи, придуманные Очиром Владимировичем, он читает 10 минут и решает тоже за 10 минут. Очир Владимирович предложил две задачи, Анджур Аркадьевич одну, а Расул Владимирович две. Какое наибольшее количество баллов Адьян успел набрать, если задачи Очира Владимировича он мог отличить, даже не читая, а отличить задачи Анджура Аркадьевича от задач Расула Владимировича может только прочитав их, при этом, читая, сначала ему попадаются задачи только Расула Владимировича. За задачи Очира Владимировича ученики получали по 4 балла, Анджура Аркадьевича - по 5 баллов, Расула Владимировича - по 8 баллов.

\newpage
\hspace{3mm}


\centerline{\textbf{\textit{Решение:}}}

Чтобы решить одну задачу РВ, нужно потратить $2+35=37$ минут. Заметим, что две задачи РВ за 1 час решить не получится (1 час = 60 минут < $37\cdot2=74$). 

Поэтому возможны два варианта: 
1. Решить одну задачу РВ, а на оставшееся время — другие задачи.
2. Вообще не решать задачи РВ и сосредоточиться на остальных.

\textbf{Первый вариант:}

Читаем одну задачу РВ, решаем её, зарабатываем 8 баллов. На это уходит 37 минут, остаётся 60-37=23 минуты. За это время мы можем сделать следующее:

1) Прочитать оставшуюся задачу РВ (2 минуты), прочитать задачу АА (2 минуты) и решить её (15 минут). Итого: 2+2+15=19 минут. Зарабатываем 5 баллов и больше ничего не успеваем сделать.

Суммарно: 8+5=13 баллов.

2) Решить две задачи ОВ за 20 минут (их читать не нужно). Зарабатываем 4+4=8 баллов и больше ничего не успеваем сделать (Этот вариант можно сделать и в другом порядке, начав с решения задачи или задач ОВ).

Суммарно: 8+8=16 баллов.

\textbf{Второй вариант:}

Задачи РВ не решаем.

Решаем две задачи ОВ за 20 минут (их читать не нужно), зарабатываем 4+4=8 баллов и остаётся 40 минут.

За это время читаем две задачи РВ (2+2=4 минуты), читаем задачу АА (2 минуты) и решаем её (15 минут). Всего тратим 2+2+2+15=21 минуту. Зарабатываем 5 баллов и больше ничего не успеваем сделать.

Суммарно: 8+5=13 баллов.

Сравнивая результаты, получаем, что максимальный балл - 16.

\textbf{Ответ: 16 баллов.}



